\documentclass[12pt,a4paper]{article}

\usepackage[utf8]{inputenc}
\usepackage[T1]{fontenc}
\usepackage[margin=2.5cm]{geometry}
\usepackage{hyperref}
\usepackage{microtype}
\usepackage{parskip}
\usepackage{lmodern}

\hypersetup{colorlinks=true, urlcolor=blue!60!black}

\pagestyle{empty}

\begin{document}

\begin{flushright}
Kundai Farai Sachikonye \\
Auf der Lichtung 19, 82178 Puchheim, Germany \\
\href{https://github.com/fullscreen-triangle}{github.com/fullscreen-triangle} \\
\texttt{kundai.sachikonye@bitspark.com} \\
(+49) 1626386344 \\[6pt]
18 May 2026
\end{flushright}

\vspace{1em}

Institut für Klinische Genetik \\
Universitätsklinikum Carl Gustav Carus Dresden \\
Ref.\ 1943

\vspace{1.5em}

\textbf{Re: Bioinformatiker*in --- Genomsequenzierung in der Onkologie (Ref.\ 1943)}

\vspace{1em}

Dear Hiring Committee,

I am applying for the bioinformatician position (Ref.\ 1943) in the Institute
for Clinical Genetics, supporting cancer genome sequencing diagnostics and
therapy identification.
The position description maps precisely to the software and pipelines I have
built independently: genome-wide variant analysis at production scale, RNA-seq
integration, machine learning applied to multidimensional cancer genomic data,
and large-scale omics pipelines with API-based database integration.

\textbf{NGS data analysis and cancer genome sequencing pipelines.}
The \href{https://github.com/fullscreen-triangle/gospel}{Gospel} framework
is a production-grade genomic analysis system processing genome-wide variant
data at $10^6$ variants per second.
The pipeline covers the full NGS analysis workflow: VCF processing and
variant filtering, structural variant detection, population frequency
annotation against gnomAD~v3.1.2, 1000~Genomes Phase~3, and UK~Biobank,
RNA-seq quantification and differential expression, functional annotation
via pathway databases (KEGG, Reactome), and protein structure analysis (RDKit).
For the oncology context specifically, Gospel implements somatic variant
analysis: distinguishing tumour-specific alterations from germline variants
using population frequency thresholds and allele frequency in tumour-normal
pairs, filtering against cancer hotspot databases, and annotating variants
with predicted functional consequences.
The pipeline is API-integrated against the major reference databases --- gnomAD,
ClinVar, COSMIC, OMIM --- using structured query patterns designed for the
kind of automated NGS data flows the position describes.
All processing is reproducible, version-controlled, and runs on distributed
computing infrastructure (Ray, Docker) scalable to clinical laboratory volumes.

\textbf{Machine learning for multidimensional cancer genomic data.}
The position lists ML applied to multidimensional biological data as a
core requirement. I have built three directly relevant systems:

The \href{https://syndrome-seven.vercel.app/tools/neoantigen}{Syndrome neoantigen predictor}
identifies which tumour somatic mutations produce immunogenic neoantigens from
the combination of MHC presentation likelihood ($\rho_\mathrm{MHC} \geq 0.72$)
and self-distinction ($\Delta\rho_\mathrm{self} > 0.15$), using spectral DFT
on three physicochemical channels with $\mathcal{O}(cL\log L)$ complexity.
For cancer genomics, this converts a tumour mutational burden VCF directly
into a ranked neoantigen list without requiring epitope database lookup ---
a pipeline step relevant to immunotherapy candidate selection in the
institute's oncology programme.

The \href{https://syndrome-seven.vercel.app/tools/resistance}{resistance predictor}
applies spectral complementarity $\rho$ to predict drug-enzyme resistance from
the enzyme sequence and drug structure: susceptibility when $\rho \geq \rho^* = 0.50$.
Validated against clinical resistance positions in TEM-1/Ampicillin,
GyrA/Ciprofloxacin, and AAC(3)-IV/Gentamicin.
For personalised cancer therapy identification, this provides a first-pass
computational screen of patient-specific enzyme variants for drug resistance
before clinical testing.

The \href{https://github.com/fullscreen-triangle/syndrome}{disease state model}
represents patient state as $\mathbf{D}\in[0,1]^8$ across eight pathophysiological
dimensions with per-dimension coherence indices; $\mathcal{O}(N\log N)$ trajectory
computation; sparse $\ell_1$ LP for minimum-intervention therapy selection.
AUC$\,=1.00$ in 25/25 loop-holonomy self-consistency validation experiments.
This provides the therapy-finding layer the position describes: given the patient's
genomic and clinical data, identify the minimum intervention set to redirect
disease trajectory to the healthy attractor.

\textbf{Pipeline maintenance, validation, and database integration.}
The position emphasises maintaining existing bioinformatic data flows and
integrating complex databases via API queries. My experience here is direct:
Gospel integrates against gnomAD, ClinVar, OMIM, COSMIC, KEGG, and Reactome
through structured REST and GraphQL API clients; Lavoisier integrates against
LIPID MAPS, LipidBlast, HMDB, and the NIST mass spectral library for automated
compound annotation. In my TUM lipidomics work I maintained and extended
multi-site federated MS analysis pipelines across distributed clinical
computing environments, including version control of pipeline configurations,
automated validation of output against reference standards, and documentation
for non-bioinformatics users.
For software validation specifically, all frameworks I have built are tested
against reference databases with documented accuracy metrics --- 98.9\,\% NIST
feature extraction for Lavoisier, AUC$\,=0.81$ vs NetMHCpan for MHC binding
prediction, $10^6$ variants/second for Gospel --- providing the kind of
quantified validation the clinical genetics context requires.

\textbf{Background.}
I hold an MSc in Computational Biology (Universität Tübingen) covering
sequence analysis, variant calling, population genetics, and bioinformatics
pipeline development; and I conducted doctoral research in Computational
Lipidomics at TUM working on large-scale MS-based molecular profiling and
multi-omics integration.
I also have direct cancer biology experience: UPLC-MS/MS shotgun proteomics
for cancer biomarker discovery at the University Medical Center Groningen (2013).
I work primarily in Python, with Rust for performance-critical components;
Java and Bash from prior work.
I am legally resident in Germany with full work authorisation and available
immediately. English is native; German is conversational.
I am willing to relocate to Dresden.

\vspace{1.5em}
Yours sincerely,

\vspace{2.5em}
\textbf{Kundai Farai Sachikonye}

\end{document}
